\documentclass[a4paper,11pt]{article}
\usepackage[utf8]{inputenc}
\usepackage[T1]{fontenc}
\usepackage[french]{babel}
\usepackage[left=1cm,right=1cm,top=.5cm,bottom=0cm]{geometry}
\usepackage{enumitem}
\usepackage{hyperref}
\usepackage{xcolor}
\usepackage{titlesec}
\usepackage{fontawesome5}
\usepackage{lmodern}
\usepackage[normalem]{ulem}

% --- Couleurs CEA ---
\definecolor{primary}{RGB}{0, 82, 147}
\definecolor{secondary}{RGB}{80, 80, 80}

% --- Configuration des liens ---
\hypersetup{
    colorlinks=true,
    linkcolor=primary,
    filecolor=primary,
    urlcolor=primary,
}

% --- Formatage des titres ---
\titleformat{\section}{\large\bfseries\color{primary}\uppercase}{}{0em}{}[\titlerule]
\titlespacing{\section}{0pt}{8pt}{5pt}

% --- Commandes personnalisées ---
\newcommand{\resumeItem}[1]{\item\small{#1}}
\newcommand{\resumeSubheading}[4]{
  \vspace{-1pt}\item
    \begin{tabular*}{0.97\textwidth}[t]{l@{\extracolsep{\fill}}r}
      \textbf{#1} & \textit{\small #2} \\
      \textit{\small #3} & \textit{\small #4} \\
    \end{tabular*}\vspace{-5pt}
}

\begin{document}

% --- EN-TÊTE ---
\begin{center}
    {\Large \textbf{Loïc Aron MBASSI EWOLO}} \\ \vspace{4pt}
    \textbf{\large Élève-Ingénieur : Deep Learning \& Mathématiques Appliquées} \\
    \textit{\small Disponible dès \textbf{Avril 2026} pour 6 mois (ou contrat de professionnalisation)} \\ \vspace{4pt}
    \small
    \faMapMarker\ Cergy, France (Mobile Île-de-France) \quad
    \faPhone\ (+33) 7 59 52 33 98 \quad
    \faEnvelope\ \href{mailto:loicaron.mbassiewolo@ensea.fr}{loicaron.mbassiewolo@ensea.fr} \\ \vspace{2pt}
    \faLinkedin\ \href{https://www.linkedin.com/in/loic-aron-mbassi-ewolo-3942a9246/}{linkedin} \quad
    \faGithub\ \href{https://github.com/Nameless0l}{github} \quad
    \faGlobe\ \href{https://mbassiloic.tech}{mbassiloic.tech}
\end{center}

\vspace{-6pt}

% --- PROFIL ---
\noindent\small{Élève-ingénieur en double diplôme ENSEA/Polytechnique, spécialisé en \textbf{Deep Learning} et \textbf{mathématiques appliquées}. Expérience de recherche au \textbf{Mila} sur la généralisation des réseaux de neurones, compétences en \textbf{PyTorch}, analyse statistique et transfert d'apprentissage. Je souhaite contribuer aux travaux sur l'IA pour les plasmas au \textbf{CEA DAM}.}

% --- FORMATION ---
\section{Formation}
\begin{itemize}[leftmargin=*, label={.}, nosep]
    \item \textbf{ENSEA (École Nationale Supérieure de l'Électronique et de ses Applications)} \hfill \textit{Cergy, France | 2025 -- Présent} \\
    \small{Ingénieur 2A, Double Diplôme - Traitement du Signal et Intelligence Artificielle.} \\
    \textit{\small{Cours : Probabilités/Statistiques, Analyse Numérique, Machine Learning, Traitement du Signal.}}
    \vspace{2pt}
    \item \textbf{École Nationale Supérieure Polytechnique de Yaoundé} \hfill \textit{Cameroun | 2021 -- 2025} \\
    \small{Diplôme d'Ingénieur de Conception (Master 2) : \textbf{Mathématiques Appliquées}, Algorithmique, Génie Logiciel.}
    \item \textbf{Université de Yaoundé I} \hfill \textit{Cameroun | 2020 -- 2022} \\
    \small{Licence Informatique \& Mathématiques : Analyse, \textbf{Algèbre Linéaire}, Probabilités, Algorithmique.}
\end{itemize}

% --- COMPÉTENCES TECHNIQUES ---
\section{Compétences Techniques}
\begin{itemize}[leftmargin=*, nosep]
    \item \small{\textbf{Deep Learning :} \textbf{PyTorch}, TensorFlow/Keras, Architectures (CNN, RNN, Transformers), Transfer Learning.}
    \item \small{\textbf{Mathématiques :} Algèbre Linéaire (SVD, décompositions), Probabilités/Statistiques, Analyse Numérique.}
    \item \small{\textbf{Data Science :} Scikit-learn, Pandas, NumPy, Analyse statistique, Visualisation (Matplotlib, Seaborn).}
    \item \small{\textbf{Programmation :} \textbf{Python} (Avancé), C/C++, MATLAB, LaTeX.}
    \item \small{\textbf{Outils :} Git, Docker, Linux, Hugging Face, Jupyter, Documentation technique.}
\end{itemize}

% --- EXPÉRIENCES DE RECHERCHE ---
\section{Expériences de Recherche}
\begin{itemize}[leftmargin=*, label={}]

    \resumeSubheading
      {Assistant de Recherche - Généralisation des Réseaux de Neurones}{Juin 2025 -- Sept. 2025}
      {MILA (Institut Québécois d'IA) - Collaboration à distance}{\textbf{Québec, Canada}}
      \begin{itemize}[leftmargin=0.4cm, nosep, label=\tiny$\bullet$]
        \item \small{Recherche sur le phénomène de « \textbf{Grokking} » (généralisation tardive après surapprentissage).}
        \item \small{Reproduction d'expériences SOTA et implémentation de pipelines sous \textbf{PyTorch}.}
        \item \small{\textbf{Analyse mathématique} de la convergence et validation par tests statistiques.}
        \item \small{Étude du comportement des réseaux de neurones profonds (MLP) sur différentes tâches.}
      \end{itemize}
    \vspace{2pt}

    \resumeSubheading
      {Projet UROP - Classification de Signaux par Deep Learning}{2024}
      {Collaboration avec mentors Google DeepMind}{Yaoundé, Cameroun}
      \begin{itemize}[leftmargin=0.4cm, nosep, label=\tiny$\bullet$]
        \item \small{Développement d'un modèle de \textbf{détection d'anomalies} sur signaux ECG.}
        \item \small{Extraction de features, filtrage numérique et prétraitement des données.}
        \item \small{Entraînement et évaluation de modèles de classification (TensorFlow/Keras).}
      \end{itemize}
    \vspace{2pt}

    \resumeSubheading
      {Stage Recherche - Machine Learning Embarqué (TinyML)}{Oct. 2023 -- Janv. 2024}
      {Laboratoire IOTA - École Polytechnique}{\textbf{Yaoundé, Cameroun}}
      \begin{itemize}[leftmargin=0.4cm, nosep, label=\tiny$\bullet$]
        \item \small{Optimisation de modèles ML pour hardware contraint (quantization, pruning).}
        \item \small{Déploiement avec TensorFlow Lite sur systèmes embarqués.}
      \end{itemize}
\end{itemize}

% --- PROJETS IA ---
\section{Projets en Intelligence Artificielle}
\begin{itemize}[leftmargin=*, label={}]

\item \textbf{Women Safe - NLP et Fine-tuning de Transformers} \hfill \textit{Compétition MLPC}
\vspace{-2pt}
    \begin{itemize}[leftmargin=0.4cm, nosep, label=\tiny$\bullet$]
        \item \small{Modèle NLP avec \textbf{BERT, GPT-2, Transformers} pour classification de textes.}
        \item \small{\textbf{Fine-tuning} de modèles pré-entraînés et \textbf{transfer learning}.}
        \item \small{Étude des biais et analyse des performances du modèle.}
    \end{itemize}
    \vspace{3pt}

\item \textbf{YembaAudio - NLP sur Langues Low-Resource} \hfill \textit{Projet Personnel}
\vspace{-2pt}
    \begin{itemize}[leftmargin=0.4cm, nosep, label=\tiny$\bullet$]
        \item \small{Adaptation de modèles \textbf{Hugging Face} à une langue peu dotée.}
        \item \small{Travail sur les données : collecte, prétraitement, augmentation.}
    \end{itemize}
    \vspace{3pt}

\item \textbf{Harmony of Languages - Vision par Ordinateur} \hfill \textit{Laboratoire IOTA}
\vspace{-2pt}
    \begin{itemize}[leftmargin=0.4cm, nosep, label=\tiny$\bullet$]
        \item \small{Création d'un \textbf{dataset custom} et augmentation de données.}
        \item \small{Entraînement de modèles de classification avec PyTorch/OpenCV.}
    \end{itemize}
    \vspace{3pt}

\item \textbf{Violence Detection - Classification Vidéo} \hfill \textit{\textbf{2\textsuperscript{ème} Place} Compétition}
\vspace{-2pt}
    \begin{itemize}[leftmargin=0.4cm, nosep, label=\tiny$\bullet$]
        \item \small{Modèle de classification temps réel sur séquences vidéo (TensorFlow, OpenCV).}
    \end{itemize}

\end{itemize}

% --- DISTINCTIONS ---
\section{Distinctions, Certifications \& Langues}
\begin{itemize}[leftmargin=*, nosep]
    \item \textbf{Hackathons :} \textbf{1\textsuperscript{ère} Place} IndabaX Africa (IA Santé), \textbf{1\textsuperscript{ère} Place} CITS National, \textbf{1\textsuperscript{ère} Place} Mountain Hub.
    \item \textbf{Leadership :} Chef Cellule IA (Polytechnique) - Animation workshops ML, Fine-tuning LLM.
    \item \textbf{Certifications :} Supervised ML (DeepLearning.AI/Stanford), Python for Data Science (IBM).
    \item \textbf{Langues :} Français (Natif), Anglais (Professionnel/Technique - Lecture articles scientifiques).
    \item \textbf{Académique :} Note Baccalauréat Physique : 19.25/20.
\end{itemize}

\end{document}
